% Figure: phases of the gait cycle, heel strike to toe off, with the
% percent-cycle annotation and the ground-reaction-force trace.
\begin{figure}[htbp]
\centering
\begin{tikzpicture}[
  axis/.style={-{Stealth}, thick, draw=engblack},
  stick/.style={thick, draw=engblack},
  label/.style={font=\scriptsize, align=center},
]

% Five stick-figure poses across the page representing gait phases.
\foreach \i/\x/\hipy/\kneey/\footx/\caption in {%
  0/0/3.4/1.7/0.2/heel strike,
  1/3.0/3.2/1.5/-0.1/loading response,
  2/6.0/3.0/1.4/-0.2/mid stance,
  3/9.0/3.2/1.5/-0.4/terminal stance,
  4/12.0/3.4/1.7/-0.6/toe off}
{
  % Head
  \draw[stick, fill=engblack!10] (\x, 4.2) circle (0.25);
  % Trunk
  \draw[stick] (\x, 3.95) -- (\x, \hipy);
  % Pelvis
  \draw[stick] (\x-0.2, \hipy) -- (\x+0.2, \hipy);
  % Stance leg
  \draw[stick, engred, very thick] (\x-0.2, \hipy) -- (\x+\footx-0.15, \kneey) -- (\x+\footx, 0.4);
  % Foot
  \draw[stick, engred, very thick] (\x+\footx-0.2, 0.4) -- (\x+\footx+0.35, 0.4);
  % Swing leg (mirror, lower stance for first/last)
  \draw[stick] (\x+0.2, \hipy) -- (\x-\footx+0.1, \kneey - 0.3) -- (\x-\footx+0.05, 0.55);
  % Caption
  \node[label] at (\x, -0.2) {\caption\\$\i\times25\,\%$};
}

% GRF curve underneath (two humps, normalised)
\begin{scope}[yshift=-3.5cm]
\draw[axis] (-0.5, 0) -- (13.5, 0) node[right, font=\small] {$t$ (\% gait cycle)};
\draw[axis] (-0.5, 0) -- (-0.5, 3.0) node[above, font=\small] {GRF $/$ BW};

\foreach \x/\xt in {0/0, 3/25, 6/50, 9/75, 12/100} {
  \draw (\x, 0) -- (\x, -0.08) node[below, font=\scriptsize] {\xt};
}
\foreach \y/\yt in {1/0.5, 2/1.0, 2.4/1.2} {
  \draw (-0.5, \y) -- (-0.6, \y) node[left, font=\scriptsize] {\yt};
}

% Double-hump curve typical of walking GRF
\draw[very thick, engblue, smooth, samples=100, domain=0:12]
  plot (\x, {2.4 * exp(-((\x - 2.0)/1.2)^2) + 2.2 * exp(-((\x - 8.0)/1.2)^2)});

\node[anchor=west, font=\small, engblue] at (2.0, 2.7) {first peak};
\node[anchor=west, font=\small, engblue] at (8.0, 2.5) {push-off};
\node[anchor=west, font=\small] at (5.0, 1.0) {valley};
\end{scope}

\end{tikzpicture}
\caption{Phases of the gait cycle (single stance leg, red) and the
typical vertical ground-reaction force trace below. The double-hump
GRF is the signature of normal walking: the first peak occurs at
loading response (the body's centre of mass falls onto the stance
leg), the valley at mid-stance (the centre of mass rises over the
stance leg, momentarily decelerating the vertical fall), and the
second peak at push-off (the ankle plantar-flexors propel the body
forward and up). Peak vertical GRF is typically $1.1$ to $1.3\,\text{BW}$
in walking; in running, the curve collapses to a single peak of
$2.5$ to $3.5\,\text{BW}$. The gait cycle conventionally runs from
heel strike (or initial contact) of one foot to the next heel strike
of the same foot.}
\label{fig:vol06ch06:gait-phases}
\end{figure}
